\section{延迟下行动的误差分析}
\label{app:analysis}

本附录推导第~\mainref{sec:formulation}~节的误差分解，把它推广到旋转、环境与分块执行，推导简单的比例关系，给出 \bench 中各项的量级，并说明其不适用的情形。这一分解促成了 \ours 的设计，它不是成功率的模型。

\subsection{平移}
\label{app:budget_translation}
设执行时刻目标的真实位置为 $\vect p(\tapp)=\vect p(\tau_t)+\vect v\Delta+\boldsymbol\rho$，策略的估计为 $\hat{\vect p}(\tapp)=\hat{\vect p}(\tau_t)+\hat{\vect v}\hat\Delta+\hat{\boldsymbol\rho}$。减去再加上 $\vect v\hat\Delta$ 得到恒等式
\begin{equation}
\begin{aligned}
\hat{\vect p}(\tapp)-\vect p(\tapp)
={}&\big[\hat{\vect p}(\tau_t)-\vect p(\tau_t)\big]+(\hat{\vect v}-\vect v)\hat\Delta\\
&+\vect v(\hat\Delta-\Delta)+(\hat{\boldsymbol\rho}-\boldsymbol\rho),
\end{aligned}
\label{eq:budget_identity}
\end{equation}
由三角不等式即得式~\mainref{eq:budget}。由于 $\vect v$ 是 $\tau_t$ 时的速度，$\boldsymbol\rho$ 恰好是匀速外推遗漏的那部分运动。在匀加速度 $\vect a$ 下，$\boldsymbol\rho=\tfrac12\vect a\Delta^2$，从历史中估计的加速度可以消除 $e_{\rm dyn}$ 的一部分；撞击、滚动接触与摩擦变化留下的余项只有动力学模型才能减小。对于学习型预测器，其估计在 $\hat{\vect v}\hat\Delta$ 与 $\hat{\boldsymbol\rho}$ 之间的划分不可辨识，因为任何预测位置都可以写成这种形式。因此我们用这一分解说明每一项需要哪些信息。我们的实验不在学习型预测器内部分离各项，而是改变作用于已知项的领先时长、速度、位姿噪声与时间，并比较接收不同信息的变体（附录~\ref{app:predictions}）。

\subsection{旋转与接触点}
\label{app:budget_rotation}
我们写 $R(\tapp)=\Exp(\boldsymbol\omega\Delta+\boldsymbol\eta)R(\tau_t)$ 与 $\hat R(\tapp)=\Exp(\hat{\boldsymbol\omega}\hat\Delta+\hat{\boldsymbol\eta})\hat R(\tau_t)$。当 $R$ 与 $\hat R$ 绕同一固定轴旋转时，指数映射可以相加，于是
\begin{equation}
\theta=\big\|\Log(\hat R(\tapp)R(\tapp)^\top)^{\vee}\big\|
\le\theta_{\rm loc}+\hat\Delta\,e_\omega+\|\boldsymbol\omega\|\,e_\tau+e_\eta ,
\label{eq:budget_rotation}
\end{equation}
其中 $\theta_{\rm loc}$ 是当前朝向误差，$e_\omega$ 与 $e_\eta$ 的定义与平移相同。对一般旋转，Baker--Campbell--Hausdorff 公式会带来角度的二阶项。物体坐标系中位于 $\vect c$、$\|\vect c\|=r$ 的接触点，其世界位置为 $\vect p+R\vect c$，因此误差为 $(\hat{\vect p}-\vect p)+(\hat RR^\top-I)R\vect c$。对转角为 $\theta$ 的旋转，$\hat RR^\top-I$ 的算子范数为 $2\sin(\theta/2)$，于是
\begin{equation}
e_{\rm c}\le\|\hat{\vect p}-\vect p\|+2r\sin(\theta/2)\le\|\hat{\vect p}-\vect p\|+r\theta ,
\label{eq:contact_app}
\end{equation}
这就是第~\mainref{sec:formulation}~节中的接触误差界。因此朝向误差以力臂为比例进入接触几何，依赖朝向的抓取或放置需要 SE(3) 的目标状态。计划中的真机球类协议使用球心平移，从输入、旋转监督及旋转误差中排除目标旋转与角运动学分量（附录~\ref{app:tracker}）。这一针对具体任务的选择不改变机器人旋转指令或仿真 SE(3) 状态，也不是一般性的对称处理结果。

\subsection{环境与间隙}
\label{app:budget_context}
令 $\mathcal P$ 为执行路径经过的环境点与时刻的集合，并设路径使机器人与被携带物体上的每一点 $\vect x$ 与\emph{预测的}环境位置 $\hat{\vect y}_t(\vect z,\tau)$（$(\vect z,\tau)\in\mathcal P$）至少相距 $m$。碰撞要求对某个 $(\vect z,\tau)\in\mathcal P$ 有 $\vect x=\vect y_t(\vect z,\tau)$，此时 $\|\hat{\vect y}_t(\vect z,\tau)-\vect y_t(\vect z,\tau)\|=\|\hat{\vect y}_t(\vect z,\tau)-\vect x\|\ge m$，因此碰撞要求路径上某处的预测误差至少为 $m$。比 $m$ 更细的环境精度因而对成功率没有一阶影响。这个界只关乎误差的大小，与采样无关。在网格上采样的流会漏掉比其间距更细的结构，并在运动边界处出错；这两种效应都进入路径误差 $e_F$，仅凭网格间距不能推断它。间隙本身也是一个设计变量：环境估计较粗时，策略可以保持更大的间隙，代价是可达性或时间。

\subsection{滚动重规划下的领先时长}
\label{app:budget_leads}
在附录~\ref{app:sweeps} 的实时执行器下，第 $t$ 次查询的总延迟 $\ell_t$ 包括附加延迟。其动作块在 $k_t\Delta_a$ 后被采用，其中 $k_t=\max(1,\lceil\ell_t/\Delta_a\rceil)$，下一次查询也在同一边界开始。设动作块长度为 $K=\cfgChunk$，在被替换或耗尽之前，到达执行阶段的指令具有领先时长 $\Delta_{t,j}=(k_t+j)\Delta_a+d_{\rm task}$，其中 $j=0,\ldots,\min(K,k_{t+1})-1$。$d_{\rm task}$ 是任务内置指令延迟：撞击后抓取为 50\,ms，其他任务为零。耗尽后的保持阶段不增加新指令。当 $k_t=k_{t+1}=k\le4$（总延迟不超过 200\,ms）时，最大领先时长为 $(2k-1)\Delta_a+d_{\rm task}$，包括任务延迟在内仍处于 \cfgMaxHorizon{} 的预测时域内。总延迟更大时，最后几条指令可能在这一时域之外执行。完整模型的动作流不接收延迟，因此不会针对特定延迟调整指令。在首个满足条件的边界重新规划可使领先时长保持较短；延迟扫描衡量更长领先时长对成功率的影响。

\subsection{比例关系}
\label{app:predictions}
这一分解给出以下参考关系。图~\ref{fig:budget_trends} 报告运动与合成噪声诊断。
\paragraph{临界领先时长。}不补偿运动的策略（$\hat{\vect v}=0$，$\hat{\boldsymbol\rho}=0$）在 $e_{\rm loc}$ 与 $\|\boldsymbol\rho\|$ 之外还要承受 $\|\vect v\|\Delta$。定位误差较小、运动平滑时，一旦 $\|\vect v\|\Delta>\varepsilon$，即超过临界领先时长 $\Delta^\star=\varepsilon/\|\vect v\|$，其接触就会失败，其成功率只通过 $\|\vect v\|\Delta/\varepsilon$ 依赖于任务、速度与延迟。表~\ref{tab:roles} 列出每个任务的 $\varepsilon$、$\|\vect v\|$ 与 $\Delta^\star$。图~\ref{fig:budget_trends}a 给出每个方法的成功率与 $\|\vect v\|\Delta/\varepsilon$ 的关系，后者在每个回合的首个指令执行时刻计算，并在已有任务与延迟条件上汇总；同时给出反应式策略的 Spearman 相关系数与拟合的 logistic 曲线。反应式策略\fillin{落在同一条曲线上}（Spearman $\rho=\fillin{X.XX}$），其成功率在 \fillin{X.X} 处减半。这一分解只给出约为 1 的阈值，而不给出曲线的形状。捕获任务在曲线周围最分散，因为其成功还取决于夹爪何时合拢。

\paragraph{位姿噪声。}由跨度为 $T$ 的 $n$ 个等间隔位姿（每个轴上带独立噪声 $\sigma$）得到的最小二乘速度估计，其标准差为 $\sigma\sqrt{12(n-1)/(n(n+1))}/T\approx\sigma\sqrt{12/n}/T$ \cite{barshalom2001}。因此在噪声达到 $\sigma^\star=\varepsilon T\sqrt n/(\Delta\sqrt{12})$ 之前，位姿历史能使 $\hat\Delta e_v$ 保持在容差以下。取 $\varepsilon=10$\,mm、$T=350$\,ms、$n=8$ 与 $\Delta=200$\,ms，得到 $\sigma^\star\approx14$\,mm。带运动先验的学习型估计器可以做得更好，因此 $\sigma^\star$ 是独立噪声下无模型线性估计对应的值。我们的鲁棒性评测采用三个合成条件：时序相关噪声、双倍幅度，以及标称噪声加连续观测缺失。它复用实时协议的无噪声参考，不估计噪声交叉点，也不声称已按真实跟踪器标定。

\paragraph{间隙。}在附录~\ref{app:budget_context} 的假设下，路径误差为 $e_F$ 的环境预测在 $m>e_F$ 时避开碰撞。这是几何上的充分条件，不是粗、细流网格在实验中等效的结论。

\begin{figure}[t]
\centering
\figplaceholder[1.5in]{\columnwidth}{图 S6——未补偿的运动与位姿噪声}{%
\textbf{(a)} x：首个指令执行时刻的 $\|\vect v\|\Delta/\varepsilon$（对数）；y：SR。每个方法 $\times$ 任务族 $\times$ 延迟一个点，复用主延迟扫描；反应式策略（灰）附拟合的 logistic 曲线与 Spearman $\rho$，标出阈值 1；补偿运动的基线与 \ours（\textcolor{predcol}{橙}）。
\textbf{(b)} x：无噪声参考、标称时序相关噪声、双倍幅度，以及标称噪声加连续观测缺失；y：SR。冻结权重并使用同一实时协议；配对区间来自匹配场景。}
\caption{\textbf{运动与位姿噪声诊断。}(a) 成功率与相对容差的未补偿运动之间的关系，比较反应式方法与补偿运动的方法。(b) 三个合成位姿噪声条件相对于无噪声参考的鲁棒性。}
\label{fig:budget_trends}
\end{figure}

\subsection{\bench 中的量级}
\label{app:budget_magnitudes}
表~\ref{tab:magnitudes} 在有代表性的 \bench 设置下（表~\ref{tab:task_sampling}）计算各项的大小。每一项都与成功条件中毫米级的容差（表~\ref{tab:tasks}）相当或更大。

\begin{table}[h]
\centering\small
\caption{领先时长为 200\,ms 时式~\mainref{eq:budget} 各项的示意量级。数值由所述量推出，并非测量值。}
\label{tab:magnitudes}
\begin{tabularx}{\columnwidth}{@{}lXr@{}}
\toprule
项 & 设置 & 大小 \\
\midrule
$\|\vect v\|\Delta$ & 传送带 0.10\,m/s & 20\,mm \\
$\hat\Delta e_v$ & 速度估计偏差 10\% & 2.8\,mm \\
$\|\vect v\|e_\tau$ & 50\,ms 时间误差（一个周期） & 7.0\,mm \\
$2r\sin(\theta/2)$ & $\omega=1.2$\,rad/s 未被预测，$r=3$\,cm & 7.2\,mm \\
$e_{\rm dyn}$ & $10^\circ$ 斜坡上的球，无加速度模型 & 24\,mm \\
$\varepsilon$ & 抓取传送带零件 & \fillin{X}\,mm \\
网格间距 & 工作空间处的 \cfgContextGrid{} 流 & \fillin{X}\,mm \\
\bottomrule
\end{tabularx}
\end{table}

斜坡一项使用纯滚动实心球的加速度 $a=\tfrac57 g\sin\alpha\approx1.2$\,m/s$^2$（$\alpha=10^\circ$），以及 $\Delta=0.2$\,s 时的 $\tfrac12a\Delta^2$。

\subsection{适用范围}
\label{app:budget_scope}
这一分解针对的是接触前运动由外部决定的刚性目标的接近过程。它不涵盖可变形或铰接的目标，因为其在 $\mathcal O$ 上的运动场不由六个数决定。它也不涵盖接触后的运动，后者取决于动作本身，需要以动作为条件的预测，这正是让世界读取动作可能起作用的情形。目标角色随时间变化的场景（例如依次操作多个物体的任务）同样不在其范围内。感知质量只通过 $e_{\rm loc}$ 与 $e_v$ 进入，我们单独测量它（表~\ref{tab:supp_tracker}）。这个界不是成功的概率，它指出哪些误差随哪些因素放大。
